\section*{ADR-007: PyO3 for Python Integration}
\addcontentsline{toc}{section}{ADR-007: PyO3 for Python Integration}

\begin{description}[leftmargin=3cm,labelwidth=2.5cm]
\item[\textbf{Status}] Accepted
\item[\textbf{Date}] April 2024
\item[\textbf{Authors}] Symphony Integration Team
\end{description}

\subsection*{Summary}

We selected PyO3 for integrating Python-based AI orchestration (Conductor) with Rust infrastructure, enabling near-zero overhead communication between the RL model and performance-critical backend components.

\subsection*{Context}

Symphony's hybrid architecture requires:
\begin{itemize}
\item Python for AI/ML capabilities and RL model implementation
\item Rust for performance-critical infrastructure
\item Low-latency communication between languages
\item Shared data structures and memory management
\item Error handling across language boundaries
\end{itemize}

\subsection*{Decision}

We implemented PyO3-based integration providing:
\begin{itemize}
\item Direct function calls from Python to Rust with FFI
\item Shared memory structures for large data transfers
\item Automatic serialization/deserialization for complex types
\item Error propagation and exception handling
\item Memory-safe resource management
\end{itemize}

\subsection*{Rationale}

PyO3 enables the performance benefits of Rust infrastructure while maintaining the AI/ML ecosystem advantages of Python for the Conductor.

\subsection*{Success Criteria}

\begin{itemize}
\item FFI overhead under 100 microseconds (Achieved)
\item Memory-safe cross-language operations (Achieved)
\item Seamless error handling (Achieved)
\end{itemize}
